
\begin{deluxetable*}{p{1.5in}p{5.0in}}
\tablecaption{Per-band quality checks. Each check is applied to each band ($g$, $r$, $i$) of every object, separately for the ZTF and PS1 sides. A side that fails any check is unusable: neither its measurement nor its detection limit enters the comparison.
\label{tab:band_filters}}
\tablehead{\colhead{Filter} & \colhead{Logic}}
\startdata
\multicolumn{2}{c}{ZTF} \\
\hline
Bright star mask & Outside the mask of every Gaia DR2 star with $G < 16$: a circle of radius $R(G) = 1630'' \times 1.396^{-G}$ and four diffraction-spike boxes $12''$ wide, of half-length $R(16) \times 1.15^{16 - G}$, at $0^\circ$, $30^\circ$, $90^\circ$ and $150^\circ$ from the detector columns \\
ZTF coverage & The position exists on a valid pixel of a ZTF reference image in the band \\
DAOPhot flags & \texttt{flags} $= 0$ for the detection (no failed or bad-data fit) \\
DAOPhot $\chi$ & \texttt{chi} $\le 6$ for the detection, the PSF-fit residual RMS over the expected (point and extended objects alike) \\
Masked pixels & $> 5$\,px from the nearest masked pixel in the band's reference image (at the detection, or at the PS1 position for a limit) \\
Kron masking & $\le 10\%$ of the ZTF Kron aperture pixels are masked (extended objects only) \\
Kron measured & Extended objects only: a ZTF detection whose PSF S/N is $\ge 5$ has a Kron measurement; without one the side cannot be judged in the photometry its morphology commits it to, and is not a non-detection either (below S/N 5 the limit stands in) \\
\hline
\multicolumn{2}{c}{PS1} \\
\hline
Bad skycells & The band's stack row is not from one of MAST's mis-registered polar skycells\footnote{\url{https://outerspace.stsci.edu/spaces/PANSTARRS/pages/298812222/Plot+of+bad+skycells+near+the+North+Pole}}\\
Catalog coverage & Without a PS1 object: the position is not in a DR2 missing-data region and $\delta > -30^\circ$\footnote{\url{https://outerspace.stsci.edu/spaces/PANSTARRS/pages/298812214/PS1+DR2+caveats}}, so the absence of an object means something \\
PS1 coverage & Without a PS1 object: some PS1 object is measured in the band within $1'$, so the band's stack reaches the position and its limit can stand in \\
Stack recovery & Without a PS1 detection in the band: the band's deep stack recovers 90\% of PS1's injected sources at some magnitude (a stack that never does, partly exposed on the survey edge or crowded in the Galactic plane, gives no limit to stand in) \\
Residual sky & Band \texttt{sky} $\le 10^{-7}$\,Jy\,arcsec$^{-2}$ \\
Stack depth & \texttt{nFrames} $\ge 4$: at least four input exposures in the stack at the position \\
Band flags & None of \texttt{DEFECT}, \texttt{PEAK\_ON\_SPIKE}, \texttt{SATSTAR} (\texttt{infoFlag}); \texttt{ON\_SPIKE}, \texttt{ON\_STARCORE}, \texttt{ON\_BURNTOOL}, \texttt{SATSTAR\_PROFILE} (\texttt{infoFlag2})\footnote{\url{https://outerspace.stsci.edu/spaces/PANSTARRS/pages/298812465/PS1+Detection+Flags}}\\
Measured & The PS1 object has a measurement in the band (forced photometry counts) \\
Kron measured & Extended objects only: as for ZTF, a PS1 measurement at PSF S/N $\ge 5$ has a Kron measurement \\
PSF coverage & \texttt{psfQfPerfect} $> 0.85$: the PSF-weighted fraction of unmasked pixels \\
PSF fit & \texttt{psfChiSq} $\le 36$, the PSF fit's reduced $\chi^2$ (the ZTF $\chi \le 6$ cut squared) \\
Primary skycell & The object has a stack row from its primary skycell (\texttt{infoFlag3} \texttt{STACK\_PRIMARY}); an object seen only in a neighbouring skycell's overlap is dropped \\
\enddata
\end{deluxetable*}

\begin{deluxetable*}{p{1.4in}p{1.5in}p{3.8in}}
\tablecaption{Per-object filters, applied after the per-band checks, in order. The first filter, ``Comparable'', sorts objects into the \texttt{in\_both}, \texttt{in\_ztf} and \texttt{in\_pstarr} catalogs. A band dropped by any check stays dropped, and an object is removed as soon as fewer than two bands remain.\label{tab:object_filters}}
\tablehead{\colhead{Filter} & \colhead{Applies to} & \colhead{Logic}}
\startdata
Comparable & All objects & A band is two-sided if both usable sides are detected ($\mathrm{S/N} \ge 5$ in the photometry the object's morphology commits it to), ZTF-only if ZTF is detected and PS1 is not, PS1-only if the reverse; a band with an unusable side, or detected by neither, compares nothing. Objects with $\ge 2$ two-sided bands go to \texttt{in\_both} (their one-sided bands are dropped); otherwise objects with $\ge 2$ ZTF-only bands and no two-sided or PS1-only band go to \texttt{in\_ztf}, and the reverse to \texttt{in\_pstarr}; everything else is removed as not comparable \\
Isolation & \texttt{in\_ztf}, \texttt{in\_pstarr} & No other ZTF or PS1 object within $3''$ \\
Sharpness & All catalogs, point sources only & $|\langle\texttt{sharp}\rangle| \le 0.5$, the $(\mathrm{S/N})^2$-weighted mean over the object's ZTF detections \\
PS1 stack detections & \texttt{in\_both}, \texttt{in\_pstarr} & \texttt{nStackDetections} $\ge 2$: PS1's own detection pass found the object in at least two stacks \\
PS1 stack quality & \texttt{in\_both}, \texttt{in\_pstarr} & Neither \texttt{BAD\_STACK} nor \texttt{SUSPECT\_STACK} set in \texttt{qualityFlag}\footnote{\url{https://outerspace.stsci.edu/spaces/PANSTARRS/pages/298812464/PS1+Object+Flags}}\\
Light not diffuse & \texttt{in\_ztf}, \texttt{in\_pstarr} & On the detecting survey's side, $(m_\mathrm{PSF} - m_\mathrm{Kron} - 0.5)/\sigma \le 3$ in all but at most one measured band: a one-sided object whose only light is an extended sheet is scattered light in the image, not a source that appeared or vanished \\
Gaia proper motion & All catalogs & Closest Gaia DR3 source within $5''$ has no proper motion, or the proper motion measurement is less than $5\sigma$ \\
Gaia parallax & All catalogs & Closest Gaia DR3 source within $5''$ has no parallax, or the parallax measurement is less than $5\sigma$ \\
Significant change & All catalogs & $|\sigma| > 3$ in at least two compared bands, with $\sigma$ the position of the band's $\Delta m = m_\mathrm{ZTF} - m_\mathrm{PS1}$ in the KDE envelope of $\Delta m$ against magnitude built from the two-sided bands of the object's morphology (point sources from PSF $\Delta m$, extended objects from Kron $\Delta m$); against a limit, only the direction the limit allows counts (brightening against a PS1 limit, fading against a ZTF limit) \\
Consistent sign & All catalogs & At least two of the bands with $|\sigma| > 3$ have $\sigma$ of the same sign \\
\enddata
\tablecomments{An object is extended if $(m_\mathrm{PSF} - m_\mathrm{Kron} - 0.05)/\sigma > 3$, combined over ZTF and PS1 by inverse variance, in at least two measured bands (in the one band if only one has a Kron magnitude); an object with no Kron magnitude anywhere is a point source. Extended objects are measured, detected and compared in Kron magnitudes on both sides (with $\mathrm{S/N} = 1.0857 / \sigma_\mathrm{Kron}$), point sources in PSF magnitudes; nothing after the morphology reads the other kind. A survey's detection limit is the magnitude at which it recovers 90\% of point sources: for ZTF from the reference image's catalog depth and pixel-noise ratio, for PS1 from its injection-recovery test of the stack the band is judged in.}
\end{deluxetable*}
